%% ============================================================
%% 分片 01：原文第 1 节「Introduction」（原文 glshort.tex 起于第 313 行）
%% 由 tools/harvmac2tex.py 从 arXiv 源码（hep-th/0604151v3）自动生成。
%% 只译英文散文；公式、编号、\label、\cite 一律不改。
%% ============================================================

\section{引言}
\label{sec:intro}

\begin{keytakeaways}
\begin{itemize}[leftmargin=1.4em,itemsep=2pt,topsep=2pt]
\item \textbf{物理图像}：把 ${\EUN}=4$ 超对称杨-米尔斯理论放在 $\Sigma\times C$ 上、令 $\Sigma$ 的尺度趋于零，低能有效理论是以 $C$ 上 Hitchin 模空间 ${\EUM}_H$ 为靶空间的二维拓扑 sigma 模型；四维的电-磁对偶（$S$-对偶）在二维约化为 Hitchin 纤维化纤维上的 $T$-对偶，即镜像对称。
\item \textbf{数学工具}：超凯勒流形 ${\EUM}_H$ 及其复结构 $I$（稳定 Higgs 丛）与 $J$（平坦 ${}^L\negthinspace G_{\Bbb{C}}$-丛）、Hitchin 纤维化与 SYZ 纤维 $T$-对偶、Fourier--Mukai 变换、Hecke 特征层与 ${\cal D}$-模、Wilson 与 't Hooft 线算子。
\item \textbf{关键结论}：GNO 双群就是 Langlands 对偶群 ${}^L\negthinspace G$（Atiyah 于 1977 年的提示）；Montonen--Olive 对偶在 ${\EUN}=4$ 理论中自然扩展为 $SL(2,\Bbb{Z})$（计入 theta 角后）；全文的六个主要想法串起「扭化 $\to$ 紧化 $\to$ 膜 $\to$ 线算子即 Hecke 算子 $\to$ ${\cal D}$-模」这条链条。
\item \textbf{边界与展望}：本文只处理几何 Langlands 对应的非分歧情形；二维共形场论在其中的作用尚未阐明；分歧情形与更一般的边界条件留待后续工作。
\end{itemize}
\end{keytakeaways}

 数域上的 Langlands 纲领 \cite{ref1} 统一了数论中许多经典与当代的成果，
是广阔的研究领域。
它对于有限域上的曲线有类似物，这方面同样有大量著名的工作 \cite{ref2,ref3}。此外，曲线的几何版本
Langlands 纲领也已得到很大的发展 
\cite{ref4,ref5,ref6,ref7,ref8}，既针对特征 $p$ 的域上的曲线，也针对通常的复黎曼面。关于一份对物理学家
相对易读、且附有大量参考文献的综述，见 \cite{ref9}。

 本文关注的是复黎曼面上的几何 Langlands 纲领。我们希望说明，这一纲领
可以如何被理解为量子场论中的一个篇章。我们不预设读者熟悉 Langlands
纲领；相反，我们预设读者熟悉超对称规范理论、电-磁对偶、sigma 模型、镜像
对称、膜以及拓扑场论这些题材。本文的主题是说明：当这些熟悉的物理要素
被用到恰好合适的问题上时，几何 Langlands 纲领会自然浮现。诸如 Hecke
特征层、${\cal D}$-模之类看似艰深的概念，会自发地出现在物理之中，并带来
关于其性质的新洞见。

 Langlands 纲领与量子场论之间存在联系的第一个线索，来自 Goddard、Nuyts
与 Olive（GNO）的工作：他们在 1976 年指出 \cite{ref10}，在规范理论中，尽管
电荷取值于规范群的权格，磁荷却取值于某个对偶群的权格。一般紧 Lie 群的
磁荷，最早由 Englert 与 Windey 分析 \cite{ref11}。GNO 的分析促成了
Montonen--Olive 电-磁对偶猜想 \cite{ref12}：基于某个给定规范群的一个具体规范
理论，等价于一个类似的、耦合常数取倒数、且规范群换为其对偶群的理论。

 对规范群 $G$ 而言，GNO 双群实际上就是 Langlands 对偶群 $^L\negthinspace G$，
后者在表述 Langlands 猜想时起着重要作用。（关于 $G$ 与 $^L\negthinspace
G$ 之间对应的若干例子，见下表。）这一点由 Atiyah 指出：1977 年末，他
向第二作者建议，Langlands 纲领与量子场论有关，并推荐了这两篇文章
\cite{ref10,ref12}。由此产生了进一步的发展 \cite{ref13}：人们认识到，
Montonen--Olive 对偶在超对称规范理论中更为自然。后来人们又认识到，
${\EUN}=4$ 超对称（即四维中可能的最大超对称）正是这一对偶的恰当语境
\cite{ref14}；而当把规范理论的 theta 角计入时，最初提出的
$\Bbb{Z}_2$ 对偶可以自然地推广为 $SL(2,\Bbb{Z})$ 
\cite{ref15,ref16}。\par\medskip \begin{table}[htbp]
\centering
\begin{tabular}{ll}
\toprule
$G$ & ${}^L G$ \\
\midrule
$U(N)$ & $U(N)$ \\
$SU(N)$ & $PSU(N)=SU(N)/\Bbb{Z}_N$ \\
$Spin(2n)$ & $SO(2n)/\Bbb{Z}_2$ \\
$Sp(n)$ & $SO(2n+1)$ \\
$Spin(2n+1)$ & $Sp(n)/\Bbb{Z}_2$ \\
$G_2$ & $G_2$ \\
$E_8$ & $E_8$ \\
\bottomrule
\end{tabular}
\caption{一个 Lie 群 $G$ 与其 Langlands 对偶（或 GNO 对偶）${}^L\negthinspace G$ 之间对应的若干例子。（原文 Table~1）}
\label{tab:gno-dual}
\end{table}

\cite{ref17}
 1990 年代初，人们猜想 Montonen--Olive 对偶可以推广到弦论
。此后，Montonen--Olive 对偶及其种种推广从许多新的视角得到
研究，并被认识到是理解场论与弦论的一个关键要素（尽管它仍然神秘）。
这些发展范围太广，无法在此综述；但那段时间的一个观察，虽然当时只是
旁支，却是本文的起点。${\EUN}=4$ 超对称杨-米尔斯理论从四维到二维的
紧化已得到研究 \cite{ref18,ref19}，并已证明低能下给出一类二维超对称
sigma 模型，其靶空间是一个超凯勒流形，即稳定 Higgs 丛的 Hitchin 模空间
${\EUM}_H$ \cite{ref20}。四维中的电-磁对偶在二维约化为 sigma 模型的
$T$-对偶。这一特定的 $T$-对偶随后在数学上被用来证明 \cite{ref21}（对
$SU(N)$ 而言）：规范群 $G$ 的 Higgs 丛模空间的 Hodge 数与
${}^L\negthinspace G$ 的相同。这一 $T$-对偶背后的几何，在 \cite{ref22}
中得到研究，随后又在 \cite{ref23} 中对任意半单 Lie 群 $G$ 得到研究。

 关于几何 Langlands 纲领与量子场论之关系的其他线索，来自相对晚近的
数学工作。Beilinson 与 Drinfeld 处理几何 Langlands 纲领的方法，建立在
$\MH$ 的量子化之上——正如他们论文的标题所示 \cite{ref5}。${\EUM}_H$
的 $T$-对偶在数学上被理解为 Fourier--Mukai 变换，它被解释为几何
Langlands 纲领的某种半经典近似。这一观点是 \cite{ref24} 一文的基础。
据了解，其他数学家（包括 Donagi 与 Pantev）也有重要的未发表贡献。
第二作者是从 D. Ben-Zvi 的一次演讲中得知 $\MH$ 的 $T$-对偶的这一
解释的；那是在 2004 年春于 IAS 举行的几何 Langlands 纲领会议上。这是
一个极其强烈的提示：必定可以用四维电-磁对偶（它导致这一特定的
$T$-对偶）以及膜（把 $T$-对偶解释为 Fourier--Mukai 变换的自然量子
场论框架）来理解几何 Langlands 纲领。这一提示正是本文的起点。

\cite{ref25}
 用尽可能简短的话来概括本文：我们将发展六个主要想法。第一个是：从
${\EUN}=4$ 超对称杨-米尔斯理论在四维中的某个扭化版本出发，可以构造
一族四维拓扑场论。在第~\ref{sec:review}~节回顾若干背景之后，我们在
第~\ref{sec:tqftfour}~节解释这一构造。扭化程序在形式上恰好类似于由
${\EUN}=2$ 超杨-米尔斯理论得到 Donaldson 理论  的构造。
第二个主要想法在第~\ref{sec:compacthitchin}~节与第~\ref{sec:tqfttwo}~节
展开：推广 \cite{ref18,ref19} 的洞见，在黎曼面 $C$ 上紧化会在二维给出
一族拓扑 sigma 模型，其靶空间为 ${\EUM}_H$，它们就是「广义 $B$-模型」。
而且，在参数的某个特殊值处，四维 $S$-对偶表现为二维镜像对称。第三个
主要想法在第~\ref{sec:loopop}~节展开：Wilson 与 't Hooft 线算子是
拓扑算子，它们以自然的方式作用在二维 sigma 模型的膜上。这里我们考虑的
是把一个膜映到另一个膜（粗略地说，把一个理论映到另一个理论）的算子，
而不是 Hilbert 空间中更熟悉的那种把一个态映到另一个态的算子。若一个
膜被 Wilson 或 't Hooft 算子映到（粗略地说）自身的倍数，我们就称它为
电特征膜或磁特征膜。$S$-对偶会自动交换电特征膜与磁特征膜。第四个主要
想法在第~\ref{sec:electricbranes}~节解释：在恰当的语境中，电特征膜与
$\pi_1(C)$ 到 Langlands 对偶群的复化 $^L\negthinspace G_{\Bbb{C}}$ 的
同态之间存在自然的对应。第五个主要想法在第~\ref{sec:thooftheckeop}~节
与第~\ref{sec:bogeqheck}~节展开：'t Hooft 算子自然地对应于几何 Hecke
算子——它们类似于几何 Langlands 纲领中的 Hecke 算子，但作用在 Higgs 丛
而非普通 $G$-丛上。还需要再一个重要的想法（在第~\ref{sec:abranes}~节
展开），才能与几何 Langlands 对应的通常表述建立联系。我们说明：由于
${\EUM}_H$ 上存在一个典范余各向同性膜，我们意义下的磁特征膜自然地与
${\cal D}$-模（微分算子层的模）相联系，后者定义在 $C$ 上全纯 $G$-丛的
模空间 ${\EUM}(G,C)$ 上。

 本文第 7、8、9、10 节以及其他各节中的部分综述性内容，主要改编自
一本即将出版的书 \cite{ref26}，该书还将包含一些额外的结果。把这些
材料纳入本文，是为了让本文更易理解。


 我们分析中的一个明显空白是：我们只考虑了几何 Langlands 对应的
非分歧情形。我们预期，把大体类似的想法用于分歧情形是可能的。第二个
主要空白是：对于二维共形场论在几何 Langlands 纲领中的用处，我们未能
有所阐明 \cite{ref5,ref9,ref27,ref28,ref29,ref30}。这些文献中的最后一篇
把共形场论用于分歧情形。我们希望，能够从本文的规范场论方法推出共形
场论方法。事实上，即使在朴素的层面上，共形场论与自守表示论之间也
存在类比 \cite{ref31}，而后者正是 Langlands 纲领的基础。最后，尽管我们
在此无可贡献，但关于 Langlands 纲领与物理之关系的又一条线索，想必来自
自守形式进入弦论的种种不同方式。若要一个极小的样本，见
\cite{ref32,ref33,ref34,ref35,ref36,ref37,ref38,ref39,ref40}。

 A.K. 感谢 D. Arinkin、R. Bezrukavnikov 与 D. Orlov 的有益讨论。特别是，
Orlov 在 2002 年关于几何 Langlands 纲领阿贝尔情形的解释，部分地促成了
\cite{ref41} 一文，该文将在第 11 节进入我们的故事。

 E.W. 感谢多年来向他解释 Langlands 纲领相关问题的众多数学家，包括
A. Beilinson、P. Deligne、V. Drinfeld 与 K. Vilonen，尤其感谢 M. F. Atiyah、
D. Ben-Zvi、R. Donagi、E. Frenkel 与 D. Kazhdan，以及最近的 M. Goresky 与
R. MacPherson。此外，T. Hausel、N. Hitchin、M. Hopkins、P. Kronheimer、
L. Jeffrey、J. Morgan、G. Moore、D. Morrison、N. Nekrasov、M. Thaddeus、
C. Vafa 与 E. J. Weinberg 等人澄清了与本文相关的若干要点；IAS 的许多
物理学家，包括 S. Hellerman、K. Intriligator、J. Maldacena、N. Seiberg 与
J. Walcher，也提出了有益的意见。


\begin{yizhu}[物理图像]
按：GNO 双群与 Langlands 对偶群之所以是同一个群，是因为二者都由同一份根数据（root data）以完全相同的方式得到：把权格与余权格互换，把单根与余单根互换。对单连通型（simply-laced）群，根与余根等长，$G$ 与 ${}^L\negthinspace G$ 因而有相同的 Lie 代数，区别只在整体形式（即中心）。这正是表中 $SU(N)$ 的对偶是 $PSU(N)=SU(N)/\Bbb{Z}_N$、$Spin(2n)$ 的对偶是 $SO(2n)/\Bbb{Z}_2$ 的原因；而 $G_2$、$E_8$ 的根数据自对偶，故它们自身对偶。
\end{yizhu}

\begin{yizhu}[易错点]
按：原文说磁荷取值于「对偶群的权格」，这句话容易与「原群的余权格」混淆，二者其实是同一个格：${}^L\negthinspace G$ 的权格按定义就是 $G$ 的余权格。因此，若只强调电荷格与磁荷格互换，两种说法等价。但要注意，电-磁对偶互换的是「格」而不只是「群」：规范群的中心与整体形式会随之改变，这正是 GNO 与 Langlands 两套语言能够对上号的关键。
\end{yizhu}

\begin{yizhu}[原文疑误]
按：原文此处记作 1976 年；GNO 一文（Goddard、Nuyts、Olive，\emph{Gauge Theories and Magnetic Charge}）正式刊于 Nuclear Phys. B125 (1977) 1--28。这里所记的年份应是其预印本流传的时间，引用时宜以期刊年份为准。
\end{yizhu}

\begin{yizhu}[物理图像]
按：这一段的物理图像值得展开。把 ${\EUN}=4$ 超对称杨-米尔斯理论放在 $\Sigma\times C$ 上，并令 $\Sigma$ 的尺度趋于零，四维理论在低能下约化为以 $C$ 上 Hitchin 模空间 ${\EUM}_H$ 为靶空间的二维 sigma 模型。${\EUM}_H$ 是超凯勒流形，带有一族复结构（$I$、$J$ 与 $K=IJ$）；四维的 $S$-对偶交换 $G\leftrightarrow{}^L\negthinspace G$，在二维则表现为 Hitchin 纤维化 ${\EUM}_H\to B$ 的纤维环面上的 $T$-对偶，也就是镜像对称的 SYZ 图景。文中提到的 Hodge 数相等，正是这一镜像对称的一个可检验推论。
\end{yizhu}

\begin{yizhu}[含糊点拨]
按：「${\EUN}=4$ 超对称杨-米尔斯理论的某个扭化版本」指的是 Kapustin--Witten 的 GL 扭化：取超对称代数中一个特定的标量超荷 $Q$（它是两个反交换超荷的线性组合，其参数取值于 $\Bbb{P}^1$），把 $Q$ 当作 BRST 算子并取 $Q$-上同调。扭化后的理论是拓扑场论，参数的选取决定二维约化后得到哪一种（广义）$B$-模型；四维 $S$-对偶作用在这些参数上，在二维就成为镜像对称。详见本文第 3、4 节。
\end{yizhu}

\begin{yizhu}[物理图像]
按：这里「特征」二字的含义要留意。Wilson 与 't Hooft 算子作用在膜的范畴上，把膜 $A$ 送到 $W(A)$；所谓「特征膜」，是指 $W(A)\cong V\otimes A$ 中的 $V$ 是一个向量空间（更一般地是一个复形或层），而不是一个数。这正是几何 Langlands 中 Hecke 特征层的范畴化版本：Hecke 算子作用在 $G$-丛模空间上的 ${\cal D}$-模范畴上，而其「特征值」是与某个局部系统相伴的向量空间。这种「算子作用在范畴上」而非「作用在态上」的现象，是本文反复使用的核心机制。
\end{yizhu}

\begin{yizhu}[含糊点拨]
按：从磁特征膜到 ${\EUM}(G,C)$ 上 ${\cal D}$-模的这一步，机制如下。${\EUM}_H$ 包含余切丛 $T^*{\EUM}_{\mathrm{st}}$ 作为 Zariski 开集（${\EUM}_{\mathrm{st}}$ 是半稳定 $G$-丛的模空间），而 $T^*{\EUM}_{\mathrm{st}}$ 上的典范余各向同性膜 $A_{\mathrm{cc}}$ 的自同态代数给出 ${\EUM}_{\mathrm{st}}$ 上的微分算子层，于是 $A\mapsto\operatorname{Hom}(A_{\mathrm{cc}},A)$ 把 A-膜送到 ${\cal D}$-模。原文此处说的是 $C$ 上全纯 $G$-丛的模空间，细节见第 6 节。
\end{yizhu}
